\chapter{Đứa Học Dốt Chưa Chắc Ngu}
\markboth{Đứa Học Dốt Chưa Chắc Ngu}{The Weak Student Is Not Necessarily Slow-Minded}

\section{Làm Chậm Nhưng Nhớ Lâu}
\begin{truyen}{Làm Chậm Nhưng Nhớ Lâu}{Slow to Finish but Able to Remember Deeply}
\chuhoa{C}{ó em làm bài rất chậm. Trong giờ, nhìn em lúc nào cũng như đi sau cả lớp một đoạn.}
\textit{(Some students work very slowly. In class, they always seem one step behind everyone else.)}

Nhiều người nhìn vậy rồi kết luận em yếu hẳn.
\textit{(Many people see that and quickly conclude the student is simply weak.)}

Nhưng vài tuần sau hỏi lại, chính em đó nhớ kỹ nhất những phần đã hiểu thật.
\textit{(But a few weeks later, that same student often remembers best what was truly understood.)}

Có những đứa trẻ không học nhanh, nhưng học chắc.
\textit{(Some children do not learn fast, but they learn solidly.)}

\end{truyen}

\begin{baihoc}
Tốc độ không nói hết chất lượng của việc học.
\textit{(Speed does not fully describe the quality of learning.)}
\end{baihoc}

\begin{ghinhoanh}
\textbf{Short English to remember:}

Speed does not fully describe the quality of learning.

\end{ghinhoanh}

\begin{tuhocnhanh}
\textbf{Gợi ý để dạy và tự nghĩ / Teaching and reflection prompts}

1. Em học sinh này yếu ở tốc độ hay yếu ở bản chất? \textit{Is this student weak in speed or in substance?}

2. Người dạy nên điều chỉnh điều gì? \textit{What should the teacher adjust?}

3. Bài học này có thể nói với học sinh ra sao? \textit{How can this lesson be shared with students?}
\end{tuhocnhanh}
\ngancach

\section{Không Nói Được, Nhưng Làm Được}
\begin{truyen}{Không Nói Được, Nhưng Làm Được}{Unable to Speak Smoothly Yet Able to Do the Work}
\chuhoa{C}{ó em lên bảng là lúng túng, trả lời miệng thì rối, càng bị nhìn càng quên.}
\textit{(Some students become confused at the board, stumble in oral answers, and forget more when being watched.)}

Nhìn vào đó, người ta dễ nghĩ em không hiểu gì.
\textit{(From that, people easily assume the student understands nothing.)}

Nhưng khi cho ngồi xuống làm chậm, từng bước một, em lại ra được bài.
\textit{(But when allowed to sit and work through things slowly, step by step, the student solves the task.)}

Có em yếu ở cách thể hiện, không phải yếu ở bên trong.
\textit{(Some students are weak in expression, not in understanding.)}

\end{truyen}

\begin{baihoc}
Đừng nhầm sự lúng túng với sự kém thông minh.
\textit{(Do not mistake awkwardness for lack of intelligence.)}
\end{baihoc}

\begin{ghinhoanh}
\textbf{Short English to remember:}

Do not mistake awkwardness for lack of intelligence.

\end{ghinhoanh}

\begin{tuhocnhanh}
\textbf{Gợi ý để dạy và tự nghĩ / Teaching and reflection prompts}

1. Em này đang yếu ở đâu? \textit{Where is this student actually weak?}

2. Thầy cô có thể tạo điều kiện gì? \textit{What condition can the teacher create?}

3. Học sinh khác cần hiểu gì về bạn mình? \textit{What should classmates understand about this student?}
\end{tuhocnhanh}
\ngancach

\section{Sợ Bảng Đen Nên Thành Ra Dở}
\begin{truyen}{Sợ Bảng Đen Nên Thành Ra Dở}{Fear of the Board Makes the Student Look Worse}
\chuhoa{N}{hiều em chỉ cần nghe gọi tên là tim đã chạy trước đầu.}
\textit{(For many students, hearing their name called is enough to make the heart run ahead of the mind.)}

Các em không sợ kiến thức bằng sợ cái cảm giác sai trước mặt bạn bè.
\textit{(They are often less afraid of knowledge than of being wrong in front of classmates.)}

Nỗi sợ đó làm một bài bình thường trở thành bài rất khó.
\textit{(That fear turns an ordinary task into a very difficult one.)}

Đôi khi điều cần dạy trước cả công thức là cảm giác an toàn.
\textit{(Sometimes what must be taught before formulas is a sense of safety.)}

\end{truyen}

\begin{baihoc}
Một lớp học tốt không chỉ dạy đúng. Nó còn phải cho phép học sinh sai mà không gãy.
\textit{(A good classroom does not only teach correctly. It also allows students to be wrong without breaking.)}
\end{baihoc}

\begin{ghinhoanh}
\textbf{Short English to remember:}

A good classroom allows students to be wrong without breaking.

\end{ghinhoanh}

\begin{tuhocnhanh}
\textbf{Gợi ý để dạy và tự nghĩ / Teaching and reflection prompts}

1. Điều gì làm em học sinh sợ nhất? \textit{What is the student most afraid of?}

2. Lớp học an toàn khác lớp học căng ở đâu? \textit{How is a safe classroom different from a tense one?}

3. Người dạy có thể đổi điều gì từ ngày mai? \textit{What can the teacher change starting tomorrow?}
\end{tuhocnhanh}
\ngancach

\section{Bị Chê Nhiều Quá Nên Đóng Cửa}
\begin{truyen}{Bị Chê Nhiều Quá Nên Đóng Cửa}{Too Much Criticism Makes a Child Shut Down}
\chuhoa{C}{ó em ban đầu còn giơ tay, còn hỏi, còn thử. Sau vài lần bị cười hoặc bị gắt, em im luôn.}
\textit{(Some students once raised their hands, asked questions, and tried. After being laughed at or scolded a few times, they went silent.)}

Người ngoài nhìn vào sẽ tưởng em lười.
\textit{(From the outside, people may assume the student is lazy.)}

Nhưng nhiều khi đó không phải lười, mà là tự vệ.
\textit{(But often it is not laziness. It is self-protection.)}

Một đứa trẻ bị tổn thương nhiều ở chỗ học rất dễ quay lưng với chuyện học.
\textit{(A child hurt too often in learning can easily turn away from learning itself.)}

\end{truyen}

\begin{baihoc}
Có những lời chê không sửa người. Nó chỉ làm người ta khép lại.
\textit{(Some criticism does not correct a person. It only closes them down.)}
\end{baihoc}

\begin{ghinhoanh}
\textbf{Short English to remember:}

Some criticism does not correct a person.

It only closes them down.

\end{ghinhoanh}

\begin{tuhocnhanh}
\textbf{Gợi ý để dạy và tự nghĩ / Teaching and reflection prompts}

1. Em này đang lười hay đang tự vệ? \textit{Is this student lazy or self-protective?}

2. Điều gì đã làm cánh cửa học tập khép lại? \textit{What closed the learning door?}

3. Mình có thể mở lại cửa đó bằng cách nào? \textit{How can that door be reopened?}
\end{tuhocnhanh}
\ngancach

\section{Chép Chậm Nhưng Nghe Rất Kỹ}
\begin{truyen}{Chép Chậm Nhưng Nghe Rất Kỹ}{Slow to Copy but Listening Closely}
\chuhoa{T}{ôi hỏi: “Sao em chép ít vậy?”}
\textit{(I asked, “Why have you copied so little?”)}

Nó nhìn tôi rồi nói nhỏ: “Em chép chậm, nên em nghe trước.”
\textit{(The student looked at me and said softly, “I copy slowly, so I listen first.”)}

Lúc đó tôi mới để ý: em không ghi nhiều, nhưng mắt không hề rời bài.
\textit{(That was when I noticed: the student wrote little, but never took their eyes off the lesson.)}

Có những em không nhanh ở tay. Các em nhanh ở chỗ bắt ý.
\textit{(Some students are not fast with their hands. They are fast at catching meaning.)}

\end{truyen}

\begin{baihoc}
Nhìn một quyển vở chưa đủ để biết một đứa trẻ có đang học thật hay không.
\textit{(One notebook is not enough to tell whether a child is truly learning.)}
\end{baihoc}

\begin{ghinhoanh}
\textbf{Short English to remember:}

One notebook is not enough to tell whether a child is truly learning.

\end{ghinhoanh}

\begin{tuhocnhanh}
\textbf{Gợi ý để dạy và tự nghĩ / Teaching and reflection prompts}

1. Em này chậm ở thao tác hay chậm ở hiểu? \textit{Is this student slow in action or in understanding?}

2. Người dạy cần quan sát thêm điều gì? \textit{What else should the teacher observe?}

3. Có bao giờ mình nhìn vở rồi kết luận quá nhanh không? \textit{Have we ever judged too quickly by a notebook?}
\end{tuhocnhanh}
\ngancach

\section{Bị Gọi Ngu Nên Càng Co Lại}
\begin{truyen}{Bị Gọi Ngu Nên Càng Co Lại}{Called Stupid and Shrinking Even More}
\chuhoa{C}{ó em nói một câu nghe rất đau: “Người ta nói em ngu riết, em cũng thấy mình ngu thiệt.”}
\textit{(One student said something painful: “People call me stupid so often that I’ve started to believe it.”)}

Một đứa trẻ nghe một nhãn quá lâu sẽ sống theo cái nhãn đó.
\textit{(A child who hears a label too long may begin living inside it.)}

Từ đó, em không dám thử nữa. Không phải vì không thể, mà vì thấy thử cũng vô ích.
\textit{(From then on, the student no longer dares to try, not because it is impossible, but because trying feels useless.)}

Nhiều học sinh yếu đi vì cách người lớn gọi tên các em.
\textit{(Many students weaken because of how adults name them.)}

\end{truyen}

\begin{baihoc}
Một lời gọi sai có thể ở lại trong đầu trẻ rất lâu.
\textit{(A wrong label can stay in a child’s mind for a very long time.)}
\end{baihoc}

\begin{ghinhoanh}
\textbf{Short English to remember:}

A wrong label can stay in a child’s mind for a very long time.

\end{ghinhoanh}

\begin{tuhocnhanh}
\textbf{Gợi ý để dạy và tự nghĩ / Teaching and reflection prompts}

1. Cái nhãn này đang làm em nhỏ đi thế nào? \textit{How is this label making the student smaller?}

2. Vì sao thử lại trở nên khó? \textit{Why does trying again become difficult?}

3. Người dạy nên thay nhãn bằng điều gì? \textit{What should replace the label?}
\end{tuhocnhanh}
\ngancach

\section{Giỏi Tay Chân Hơn Giỏi Giấy}
\begin{truyen}{Giỏi Tay Chân Hơn Giỏi Giấy}{Better with Hands than with Paper}
\chuhoa{C}{ó em làm bài lý thuyết chật vật, nhưng hễ giao việc thực hành là rất nhanh.}
\textit{(Some students struggle in written theory, yet move quickly when given practical work.)}

Tôi bảo: “Em làm cái này được đó.”
\textit{(I said, “You’re good at this.”)}

Nó cười một cái rất lạ, kiểu cười của một đứa ít khi được khen đúng chỗ.
\textit{(The student smiled in a strange way, the smile of someone rarely praised in the right place.)}

Không phải đứa trẻ nào cũng sáng lên ở cùng một dạng bài.
\textit{(Not every child lights up in the same kind of task.)}

\end{truyen}

\begin{baihoc}
Giá trị của học sinh nhiều khi nằm ở chỗ bài kiểm tra chưa chạm tới.
\textit{(A student’s value often lives in places ordinary tests never reach.)}
\end{baihoc}

\begin{ghinhoanh}
\textbf{Short English to remember:}

A student’s value often lives in places ordinary tests never reach.

\end{ghinhoanh}

\begin{tuhocnhanh}
\textbf{Gợi ý để dạy và tự nghĩ / Teaching and reflection prompts}

1. Em này đang sáng ở dạng năng lực nào? \textit{In what kind of ability does this student shine?}

2. Vì sao lời khen đúng chỗ quan trọng? \textit{Why does accurate praise matter so much?}

3. Mình có đang chỉ nhìn một kiểu giỏi không? \textit{Are we only recognizing one kind of intelligence?}
\end{tuhocnhanh}
\ngancach

\section{Một Lần Được Tin Là Khác Hẳn}
\begin{truyen}{Một Lần Được Tin Là Khác Hẳn}{One Moment of Trust Changes the Student}
\chuhoa{T}{ôi giao cho một em học yếu việc phát bài. Có đứa nhìn sang như không hiểu vì sao tôi lại chọn em đó.}
\textit{(I asked a weak student to hand out the worksheets. Some classmates looked surprised that I chose that child.)}

Em làm rất cẩn thận, đi rất chậm, nhưng mặt rất sáng.
\textit{(The student worked carefully, slowly, but with a bright face.)}

Có những đứa trẻ học dở một thời gian dài chỉ vì ít khi được giao một điều gì chứng tỏ rằng người lớn vẫn tin mình.
\textit{(Some children stay weak for a long time simply because they are rarely given anything that says an adult still trusts them.)}

Được tin đôi khi là bài học đầu tiên trước khi học lại kiến thức.
\textit{(Being trusted can be the first lesson before relearning knowledge.)}

\end{truyen}

\begin{baihoc}
Niềm tin đúng lúc có thể mở đầu cho một sự thay đổi rất thật.
\textit{(Timely trust can begin a very real change.)}
\end{baihoc}

\begin{ghinhoanh}
\textbf{Short English to remember:}

Timely trust can begin a very real change.

\end{ghinhoanh}

\begin{tuhocnhanh}
\textbf{Gợi ý để dạy và tự nghĩ / Teaching and reflection prompts}

1. Vì sao được tin lại quan trọng với học sinh yếu? \textit{Why does trust matter so much for weak students?}

2. Người dạy có thể giao việc gì để tạo niềm tin? \textit{What task can a teacher give to build trust?}

3. Mình đang giữ niềm tin ở học sinh nào quá ít? \textit{Which student are we trusting too little right now?}
\end{tuhocnhanh}
